% 13_bh_interior_regime_iv.tex
% Regime IV: The Black Hole Interior as the DC Operating Point Past Breakdown
\chapter{Regime IV: The Black Hole Interior}
\label{ch:bh_interior}


\begin{objectivebox}
\begin{itemize}
    \item Model the black hole interior as Regime~IV: the DC operating point past dielectric breakdown with finite saturation impedance.
    \item Understand why no singularity forms and how the information paradox is eliminated by preserving finite internal degrees of freedom.
\end{itemize}
\end{objectivebox}

This chapter characterizes \textbf{Regime~IV (Ruptured Topology)}: the
saturated black-hole interior. There is a fundamental structural distinction
between the black hole interior and particle confinement, despite both being
instances of $S(A/A_c) = 0$.

\section{Motivation: The Missing Regime}

The AVE framework defines four universal regimes via $S(r) = \sqrt{1 - r^2}$.
Prior work validates Regimes~I--III:
\begin{itemize}
    \item \textbf{Regime I:} WD redshift ($\varepsilon_{11} \sim 10^{-3}$),
          LIGO ($h \sim 10^{-21}$), periodic table.
    \item \textbf{Regime II:} BCS at $T/T_c \approx 0.5$, Asymmetric Active Capacitor at 30~kV.
    \item \textbf{Regime III:} BCS near $T_c$, Asymmetric Active Capacitor at 43~kV.
\end{itemize}
Regime~IV ($S = 0$, topology destroyed) has \textit{no engine validation}.
The regime map claims neutron stars at $\varepsilon_{11} = 1.46$
and black holes at $\varepsilon_{11} = 3.5$ both lie in Regime~IV,
but no computation profiles the interior properties.

\section{Step 1: LC Analogs}

Under Symmetric Gravity (Axiom~3), both $\mu$ and $\varepsilon$ scale by $n(r)$:
\begin{equation}
    \mu'(r) = \mu_0 \cdot n(r), \qquad \varepsilon'(r) = \varepsilon_0 \cdot n(r)
\end{equation}
The impedance is therefore invariant:
\begin{equation}
    Z(r) = \sqrt{\frac{\mu'(r)}{\varepsilon'(r)}} = \sqrt{\frac{\mu_0}{\varepsilon_0}} = Z_0
    \qquad \text{(for all $r$, including $r < r_{sat}$)}
\end{equation}
This means there is \textbf{no EM impedance mismatch} at any gravitational
boundary: in the EM/transverse channel $Z_{EM} = Z_0$ and the reflection
coefficient $\Gamma_{EM} = (Z_2 - Z_1)/(Z_2 + Z_1) = 0$ everywhere. The EM
refractive index $n(r) - 1 = 2GM/(c^2 r)$ reaches unity at the Schwarzschild
radius $r_s = 2GM/c^2$ (where light sees the GR horizon, $n \to \infty$); this
is the EM horizon, not a $\Gamma = -1$ reflector. The shear and bulk channels
are treated separately: at $r_{sat} = 7GM/c^2 = 3.5\,r_s$ the shear modulus
collapses ($G_{shear} \to 0$, $Z_{shear}, Z_{bulk} \to 0$, $\Gamma = -1$).

\begin{resultbox}{The DC Operating Point}
In electrical engineering terms, gravity is a matched transmission line.
The BH interior is the region past dielectric breakdown ($V > V_{BR}$):
the DC operating point has exceeded the yield voltage.
No small-signal (AC) model exists. Transverse waves (shear/GW)
cannot propagate because the shear modulus is zero.
\end{resultbox}

\section{Step 2: Strain and Regime Profile}

For a 10~$M_\odot$ black hole:
\begin{center}
\begin{tabularx}{\textwidth}{@{}cccccX@{}}
\toprule
$r/r_{sat}$ & $r$ [km] & $\varepsilon_{11}$ & Regime & $S$ & $c_{eff}/c$ \\
\midrule
10.0  & 1034 & 0.100 & I  & 0.995 & 0.997 \\
3.0   & 310  & 0.333 & II & 0.943 & 0.971 \\
1.5   & 155  & 0.667 & II & 0.745 & 0.863 \\
1.2   & 124  & 0.833 & II & 0.553 & 0.744 \\
1.0   & 103  & 1.000 & IV boundary & 0 & 0 \\
0.5   & 52   & 2.000 & IV & 0 & 0 \\
0.286 & 30   & 3.500 & IV & 0 & 0 \\
\bottomrule
\end{tabularx}
\end{center}

\noindent Key checkpoints:
\begin{itemize}
    \item $r_{sat} = 7GM/c^2 = 103.4$~km, $r_s = 2GM/c^2 = 29.5$~km
    \item $r_{sat}/r_s = 7/2 = 3.50$ (exact, from $\varepsilon_{11} = 7GM/(c^2 r)$)
    \item $\varepsilon_{11}(r_s) = 7/2 = 3.50$ (exact)
\end{itemize}

\section{Step 3: The 0$\cdot\infty$ Limit}

At $r \to 0$:
\begin{equation}
    n(r) = 1 + \frac{2GM}{c^2 r} \to \infty, \qquad S(r) = 0
\end{equation}
Therefore:
\begin{equation}
    \mu_{eff} = \mu_0 \cdot n \cdot S = \mu_0 \cdot \infty \cdot 0
    = \text{indeterminate}
\end{equation}
This is not a pathology: in the ruptured interior, the lattice topology is destroyed.
The constitutive parameters $\varepsilon$ and $\mu$ describe the elastic compliance
of a structure that no longer exists. The $0 \cdot \infty$ singularity
is the signature that the LC model breaks down---it is analogous
to querying the AC impedance of a burned-out semiconductor device.

\section{Step 4: The Regime IV Isomorphism}

All four Regime~IV systems use $S(A/A_c) = 0$, but they differ
in \textit{which} constitutive parameter saturates:

\begin{resultbox}{Regime IV Isomorphism Table}
\begin{center}
\begin{tabularx}{\textwidth}{@{}lll>{\centering\arraybackslash}X>{\centering\arraybackslash}X>{\centering\arraybackslash}XX@{}}
\toprule
\textbf{System} & \textbf{Saturates} & \textbf{Sym} & $Z$ & $\Gamma$ & \textbf{Interior} & \textbf{EE Analog} \\
\midrule
Electron          & $\mu$ (inductive) & Asym & $\to 0$   & $-1$ & Standing wave & LC tank \\
Superconductor    & $\mu_{eff}$       & Asym & $\to 0$   & $-1$ & Standing wave & Shorted $L$ \\
Nucleus (Pauli)   & $U(r)$ at $d_{sat}$ & Asym & $\to\infty$ & $+1$ & Repulsive wall & Open circuit \\
BH interior       & $G_{shear}$       & \textbf{Sym}  & $Z_{EM}{=}Z_0$; $Z_{shear},Z_{bulk}{\to}0$   & $\Gamma_{EM}{=}0$; $\Gamma_{shear}{=}\Gamma_{bulk}{=}{-}1$  & EM-transparent; shear/bulk-reflecting & Past $V_{BR}$ \\
\bottomrule
\end{tabularx}
\end{center}
\end{resultbox}

\noindent
\textbf{The BH is the ONLY Regime~IV system with no EM-channel impedance
mismatch} ($\Gamma_{EM} = 0$); its shear and bulk channels still reflect
($\Gamma_{shear} = \Gamma_{bulk} = -1$ at $r_{sat}$).
The electron is a \textit{knot} (topology wound up);
the BH is a \textit{hole} (topology torn apart).
Both use the same saturation operator, but the symmetry of saturation
determines whether the interior is a resonant cavity (EM-reflecting) or an
EM-transparent sink that still reflects shear/GW modes.

\section{Step 5: Characteristic Scales}

\begin{center}
\begin{tabularx}{\textwidth}{@{}lccX@{}}
\toprule
\textbf{Object} & $r_{sat}$ or $r_{conf}$ & $\varepsilon_{11}$ at surface & Regime \\
\midrule
10 $M_\odot$ BH     & 103.4 km  & 1.0 (at $r_{sat}$) & IV boundary \\
62 $M_\odot$ BH     & 641.0 km  & 1.0 (at $r_{sat}$) & IV boundary \\
Sgr A* (4.3M $M_\odot$) & 44,460 Mm & 1.0 (at $r_{sat}$) & IV boundary \\
Electron ($0_1$ unknot, real-space) & $\sim$0.39 fm & $\to 1$ from below & IV boundary \\
Proton ($6^3_2$ Borromean linkage, real-space)   & $\sim$0.87 fm & $\to 1$ from below & IV boundary \\
\bottomrule
\end{tabularx}
\end{center}

\noindent
Despite spanning 23 orders of magnitude in size, all five objects sit
at the same \textit{universal boundary}: $S = 0$, $r = A/A_c = 1$.
The physics is scale-invariant; only the physical meaning of the
control parameter changes.

\section{Step 6: Testability and Predictions}

\begin{enumerate}
    \item \textbf{No information paradox (EM/thermodynamic channel).}
    In the EM channel the boundary is matched ($\Gamma_{EM} = 0$) and the
    interior is a dissipative sink: EM/topological information is erased at
    the ruptured boundary (the thermodynamic floor), not ``stored'' in a
    unitary cavity. This resolves the unitarity debate in the EM channel:
    the question ``where does information go?'' has the same answer
    as ``where does charge go when a capacitor breaks down?'' The shear and
    bulk channels reflect rather than absorb ($\Gamma_{shear} = \Gamma_{bulk}
    = -1$ at $r_{sat}$), which is the basis of the GW-echo prediction.

    \item \textbf{Hawking radiation = saturation-boundary noise leakage.}
    Hawking radiation is the classical thermodynamic leakage of ambient
    lattice (Nyquist) noise through the imperfect saturation boundary at
    $r \approx r_{sat}$, with $T_H = \hbar c^3/(8\pi G M k_B)$ --- not
    quantum tunnelling. It is the gravitational analogue of
    saturation-boundary emission at $E = E_{yield}$ (same $S \to 0$
    boundary, different constitutive sector).

    \item \textbf{BH merger ringdown (validated).}
    The collision of two ruptured regions creates a new $r_{sat}$ boundary.
    The ringdown IS the exterior lattice equilibrating to the new
    saturation shell. Already validated to 1.7\% via the 5-step
    eigenvalue method ($\omega \cdot M = 0.3673$ vs.\ GR exact $0.3737$).
\end{enumerate}

\section{Conclusions}

\begin{enumerate}
    \item The BH interior is Regime~IV of the universal saturation operator:
    $S(\varepsilon_{11}) = 0$ for $r < r_{sat} = 7GM/c^2$.

    \item Unlike particle confinement (asymmetric saturation, $\Gamma = -1$),
    the BH is \textbf{symmetric in the EM channel}: $Z_{EM} = Z_0$ everywhere,
    $\Gamma_{EM} = 0$ --- EM-transparent, a dissipative sink in that channel.
    The shear and bulk channels still reflect at $r_{sat}$ ($Z_{shear},
    Z_{bulk} \to 0$, $\Gamma_{shear} = \Gamma_{bulk} = -1$), so the interior
    is a perfect reflector for shear/GW modes (the GW-echo falsifier).

    \item The $0 \cdot \infty$ singularity at $r \to 0$ is the signature
    that the LC model breaks down in the ruptured interior---analogous
    to querying the AC model of a destroyed semiconductor device.

    \item In EE terms: the BH interior is the DC operating point
    past dielectric breakdown. The ``small-signal model'' (GW propagation)
    does not exist because the ``device'' (lattice topology) is destroyed.

    \item The exterior supports shear waves; the QNM eigenfrequencies
    arise from $r_{eff} = r_{sat}/(1+\nu_{vac})$, the same formula
    used for WD surface modes and validated to 1.7\% against GR.
\end{enumerate}

\section{The substrate-observability horizon — same mechanism at every scale}
\label{sec:substrate_observability_horizon}

\textbf{Cross-cutting note.}
The BH interior's epistemic horizon --- no observer outside the event horizon
can see the interior matter-history; only the three integrated invariants
$M, Q, J$ are externally measurable --- is the SAME mechanism that walls
off our cosmological initial conditions from internal observers. The
substrate-observability rule (Common Foreword $\S$``Three Boundary
Observables and the Substrate-Observability Rule'') applies fractally at
every scale where a $\Gamma = -1$ boundary forms:

\begin{itemize}
\item \textbf{BH exterior $\to$ BH interior:} outside observers measure
  $M, Q, J$ at the event horizon. The matter-history that formed the BH
  is fundamentally inaccessible from outside. No observer position outside
  the horizon can resolve interior structure.

\item \textbf{Our cosmic interior $\to$ pre-genesis state:} we sit inside
  our cosmic $\Gamma = -1$ boundary (cosmic horizon $R_H \approx 10^{26}$~m,
  identified per Vol~3 Ch~4 canonical as our parent BH's Schwarzschild
  radius). From inside, we measure the cosmic boundary's three observables
  $\mathcal{M}_{\text{cosmic}}, \mathcal{Q}_{\text{cosmic}}, \mathcal{J}_{\text{cosmic}}$
  via local-physics consequences (CMB axis anomalies, large-scale-structure
  rotation correlations, Hubble flow anisotropy, frame-dragging at cosmic
  scale). The \emph{specific cosmic parameters} ($M_{\text{parent BH}}$,
  $J_{\text{parent BH}}$) that set those observables at lattice genesis
  are fundamentally inaccessible from inside.
  \emph{The mechanism itself is not inaccessible}
  per the A-034 catalog (Vol~3 Ch~4 \S\ref{sec:tki_strain_snap}): it is the
  universal saturation-kernel strain-snap mechanism, directly observable at
  four other scales (atomic dielectric breakdown;
  planetary geomagnetic reversal in Ch~13; stellar solar flare in Ch~14;
  BH ring-down QNM in Ch~15).

\item \textbf{Same epistemic horizon, different scales (A-034 catalog).}
  The crystallization phase transition that produced our substrate IS the
  wall --- no observer exists within the resulting substrate from a position
  outside it. This is structurally identical to a Kerr BH astronomer's
  relationship to the interior: they characterize the boundary, not the
  specific parameters that set it. The MECHANISM CLASS (Axiom~4 strain-snap)
  IS observable: we see it at four smaller scales, giving cross-scale
  validation paths.
\end{itemize}

This cross-cutting framing connects Ch~\ref{ch:bh_interior} to Vol~3 Ch~4
(Generative Cosmology cosmic-boundary canonical + \S\ref{sec:tki_strain_snap}
TKI Strain-Snap) and Backmatter Ch~\ref{app:universal_saturation_kernel}
(Universal Saturation-Kernel Catalog with all 26 cross-scale instances).
The substrate-observability rule is canonical at every scale AND canonical
when applied to ourselves --- the framework's epistemic horizon for
\emph{specific cosmic parameters}, not for the mechanism class.

\begin{summarybox}
\begin{itemize}
    \item The black hole interior is modelled as Regime~IV: the DC operating point past dielectric breakdown where the vacuum lattice is fully saturated ($r_{\text{eff}} = r_{\text{sat}}/(1 + \nu_{\text{vac}})$).
    \item No singularity forms; the interior approaches a finite-impedance saturation state that smoothly connects to the exterior Regime~III Schwarzschild solution.
    \item Predictions match GR exterior observables to $1.7\%$ while eliminating the information paradox by preserving finite internal degrees of freedom.
    \item \textbf{Same-epistemic-horizon framing:} the BH interior's
          inaccessibility from outside is the same substrate-observability
          mechanism that walls off our cosmological initial conditions
          from internal observers. We characterize boundaries (3 observables),
          not what set them.
\end{itemize}
\end{summarybox}

\begin{exercisebox}
\begin{enumerate}
    \item Explain why the AVE framework predicts no singularity at the centre of a black hole, using the Axiom~4 saturation limit.
    \item Derive $r_{\text{eff}} = r_{\text{sat}}/(1 + \nu_{\text{vac}})$ and explain its physical interpretation as the minimum resolvable length scale inside the event horizon.
\end{enumerate}
\end{exercisebox}

